\chapter{Предисловие}

\epigraph{Physics is becoming so unbelievably complex that it is taking longer
and longer to train a physicist. It is taking so long, in fact, to train a 
physicist to the place where he understands the nature of physical problems that 
he is already too old to solve them.}{Eugene Paul Wigner}

Я физик.Физик по образованию, образу мысли и, в целом, по профессии.
 Моя профессиональная жизнь разделилась на две части. После окончания Новосибирского 
университета в 1977 году, я в течение 15 лет занимался теоретической физикой 
(в основном квантовой теорией поля). В 1992 году я репатриировался в Израиль  
(о чем, кстати, ни разу не пожалел), и моя академическая карьера закончилась.
Я начал работать в промышленности. Занимался разработкой алгоритмов и програм, 
моделированием и анализом данных, вычислениями и численными экспериментами.

Я называл эту свою деятельность промышленной физикой, так как во главу угла
ставил понимание сути процесса или явления в противоположность чисто статистическим
подходам типа нейронных сетей, машинного обучения и прочим разновидностям
искусственного интелекта.

 Сейчас я на пенсии и решил поделиться своими мыслями о теоретической физике.

